\section{Универсальные семейства}

\subsection{Домашнее задание}
\begin{enumerate}
  \item
    Пусть $A$ --- непустое конечное множество и $\mathcal{A} = \{ f : A \rightarrow \F_2^n \}$ --- 2-независимое семейство хеш-функций на нем.
    Обозначим за $L_k$ множество списков элементов из $A$ длины \textbf{ровно} $k$, за $L_{\le k}$ --- длины \textbf{не более} $k$ и за $L$ --- списков любой длины.
    \begin{enumerate}
      \item
        Фиксируем $k$ и введем семейство хеш-функций $\mathcal{H} = \{ h : L_k \rightarrow \F_2^n \}$ следующим образом:
        пусть $c \in \F_2^n$ и $f_0, f_1, \dots, f_{k-1} \in \mathcal{A}$ семплированы случайно независимо.
        Тогда: $$h([ l_0, l_1, \dots, l_{k-1} ]) = c \oplus f_0(l_0) \oplus f_1(l_1) \oplus \dots \oplus f_{k-1}(l_{k-1})$$
        Покажите, что полученное семейство 2-независимое. \\
        \textit{Подсказка: используйте результат с практики и индукцию по $k$.}

      \item
        Теперь аналогично определенное семейство хеш-функций $\mathcal{H} = \{ h : L_{\le k} \rightarrow \F_2^n \}$:
        независимо семплируем $c \in \F_2^n$ и $f_0, f_1, \dots, f_{k-1} \in \mathcal{A}$, и:
        $$h(l) = c \oplus f_0(l_0) \oplus f_1(l_1) \oplus \dots \oplus f_{len(l) - 1}(l_{len(l) - 1})$$
        (тут $len(l) \le k$).
        Покажите, что полученное семейство также 2-независимое. \\
        \textit{Подсказка: не забудте про пустой список.}

      \item
        Придумайте 2-независимое семейство $\{ L \rightarrow \F_2^n \}$.
        Если оно получится бесконечным, также придумайте практический способ как им пользоваться.
    \end{enumerate}

  \item
	Равномерной называется хеш-функция, у которой у каждого значения одинаковое количество прообразов.
	Пусть есть пара множеств $A$, $B$ и две равномерные хеш-функции $f : A \rightarrow [0 \dots n-1]$ и
	$g : B \rightarrow [0 \dots n-1]$, где $n$ --- натуральное число.
	Постройте равномерную хеш-функцию с областью значений $[0 \dots n-1]$ для:
	\begin{enumerate}
	  \item Упорядоченных пар $(A, B)$,
	  \item Неупорядоченных пар $\{ A, A \}$, если $n$ --- нечётное.
	  \item Покажите, что при четном $n$ соответствующей хеш-функции для $\{ A, A \}$ может не существовать.
	\end{enumerate}

  \item
    Есть семейство хеш-функций $\mathcal{H} = \{ h : K \rightarrow Y \}$ и подмножество
    $X \subseteq K$, $|X| > 1$. Покажите, что, если $\mathcal{H}$ универсальное
    и $|Y| = c |X|^2$, то
    $$\Prb{h \in \mathcal{H}}\left[ \; |h(X)| = |X| \; \right] \ge 1 - \frac{1}{2 c}$$

    Подсказка: вам может пригодиться \href{https://ru.wikipedia.org/wiki/%D0%9D%D0%B5%D1%80%D0%B0%D0%B2%D0%B5%D0%BD%D1%81%D1%82%D0%B2%D0%BE_%D0%9C%D0%B0%D1%80%D0%BA%D0%BE%D0%B2%D0%B0}{неравенство Маркова}.

\end{enumerate}
